
\begin{proposition}[Rational normalization bounds with the exact constant-sheaf map]\label{cdp:st:A.5.16}
Let $f:Y\to B$ be a proper holomorphic map from a smooth compact
complex threefold with $b_2(Y)=0$, with connected fibres and smooth
general fibre a two-torus. Let $mS$ be a full fibre with irreducible
nonnormal reduced support. Write $\eta:T\to S$ for normalization,
$\pi:Z\to T$ for its minimal resolution, and $\Gamma\subset S$
for the conductor-image scheme. Suppose that $Z$ is rational.
Put $Q=\eta_*\OO_T/\OO_S$ and $e=h^0(S,\omega_S)$.
Then $e\in\{0,1\}$, and
\[
\begin{gathered}
 R^1\pi_*\OO_Z=0,\qquad
 H^1(T;\mathbb C)=H^1(T,\OO_T)=H^2(T,\OO_T)=0,\\
 H^0(S,Q)\xrightarrow{\ \simeq\ }H^1(S,\OO_S),\qquad
 h^0(S,Q)=1+e,\quad h^1(S,Q)=e,\quad
 \chi(\Gamma,\OO_\Gamma)=-1.
\end{gathered}
\tag{A.RT.1}\label{cdp:eq:A.RT.1}
\]
In particular $T$ has rational singularities. The coefficient map
\[
 H^1(S;\mathbb C)\longrightarrow H^1(S,\OO_S)
\tag{A.RT.2}\label{cdp:eq:A.RT.2}
\]
is injective, so $b_1(S)\le1+e\le2$.

There is an effective integral divisor $D$ on $Z$ with
$\OO_Z(D)\simeq K_Z^{-1}$, whose strict-transform part is the
normalization conductor. It satisfies
\[
 h^0(D,\OO_D)=h^1(D,\OO_D)=1,
 \quad D_{\rm red}\text{ connected},
 \quad p_a(D_{\rm red})\le1.
\tag{A.RT.3}\label{cdp:eq:A.RT.3}
\]
The conductor image is connected,
$p_a(\Gamma_{\rm red})\le2$, and
$e(\Gamma_{\rm red})\ge0$. Set-theoretically
$D_{\rm red}=\pi^{-1}(\eta^{-1}(\Gamma_{\rm red}))$.
Exceptional curves over points outside the conductor have
coefficient zero in $D$; they are not to be counted as conductor
branches. With $E_p=(\pi^{-1}(p))_{\rm red}$ the exact Euler
formula is
\[
 e(S)=e(Z)-e(D_{\rm red})+e(\Gamma_{\rm red})
       -\sum_{p\in\operatorname{Sing}T\setminus\eta^{-1}(\Gamma)}
                  \bigl(e(E_p)-1\bigr).
\tag{A.RT.4}\label{cdp:eq:A.RT.4}
\]
In the torus smoothing of the statement, \ref{cdp:st:A.5.17} excludes
every point in this sum. Hence $T$ is smooth outside the
normalization conductor, and the sum is zero.
\end{proposition}
\begin{proof}
The canonical evaluation argument in \ref{cdp:st:A.5.12} gives
$\omega_S^m\simeq\OO_S$ and $\chi(S,\OO_S)=0$.
Thus $e=0$ or $1$, with $e=1$ precisely when
$\omega_S\simeq\OO_S$, and
$h^1(S,\OO_S)=1+e$.

We first prove a useful fact without the rationality assumption.
Finite duality identifies the normal dualizing sheaf with
\[
 \omega_T=\mathfrak c_T\otimes\eta^*\omega_S
       \ \subset\ \eta^*\omega_S,
\tag{A.RT.5}\label{cdp:eq:A.RT.5}
\]
where $\mathfrak c_T$ is the conductor ideal on $T$.
A nonzero global section of $\omega_T$ is therefore a nonzero
section of the torsion line $\eta^*\omega_S$ vanishing on its
nonempty conductor. Its $m$th tensor power would be a nonzero
holomorphic function on compact connected $T$ vanishing there,
which is impossible. Consequently
$H^2(T,\OO_T)=H^0(T,\omega_T)^\vee=0$.

For the proper resolution $\pi$, normality gives
$\pi_*\OO_Z=\OO_T$, and its fibres have dimension at most one.
The low-degree Leray sequence is
\[
\begin{split}
0\longrightarrow H^1(T,\OO_T)\longrightarrow H^1(Z,\OO_Z)
 \longrightarrow H^0(T,R^1\pi_*\OO_Z)\\
 \longrightarrow H^2(T,\OO_T)\longrightarrow H^2(Z,\OO_Z)
 \longrightarrow0.
\end{split}
\tag{A.RT.6}\label{cdp:eq:A.RT.6}
\]
Rationality of $Z$ makes both its positive-degree groups zero.
Together with \eqref{cdp:eq:A.RT.5}, this gives $H^1(T,\OO_T)=0$ and
$R^1\pi_*\OO_Z=0$, since this last sheaf has finite support.
The fibres of $\pi$ are connected. Proper base change for the
constant sheaf gives $R^0\pi_*\mathbb C_Z=\mathbb C_T$.
The degree-one topological Leray edge map is consequently an
injection $H^1(T;\mathbb C)\hookrightarrow H^1(Z;\mathbb C)=0$.

In the normalization exact sequence
\[
 0\longrightarrow\OO_S\longrightarrow\eta_*\OO_T
       \longrightarrow Q\longrightarrow0,
\tag{A.RT.7}\label{cdp:eq:A.RT.7}
\]
the map on $H^0$ is the identity on constants. The preceding
vanishings therefore make its connecting map
$H^0(Q)\to H^1(\OO_S)$ an isomorphism and its next connecting
map $H^1(Q)\to H^2(\OO_S)$ an isomorphism. These are the
maps and dimensions in \eqref{cdp:eq:A.RT.1}. In particular $\chi(Q)=1$.
The exact conductor-duality isomorphism proved in \ref{cdp:st:A.5.9} is
\[
 Q\simeq\omega_\Gamma\otimes\omega_S^{-1}|_\Gamma.
\tag{A.RT.8}\label{cdp:eq:A.RT.8}
\]
The conductor image is a proper Cohen--Macaulay curve, including
any nilpotent structure. Serre duality on it gives
$\chi(Q)=-\chi(\omega_S|_\Gamma)$. A torsion line bundle has
degree zero on every reduced component, and therefore also has
total degree zero on this possibly nonreduced curve. Curve
Riemann--Roch gives
$\chi(\omega_S|_\Gamma)=\chi(\OO_\Gamma)$.
This proves $\chi(\OO_\Gamma)=-1$.

Here is the asserted injectivity of the complex coefficient map.
In the category of sheaves of complex vector spaces set
\[
 \mathcal V=\operatorname{coker}
     (\mathbb C_S\longrightarrow\eta_*\mathbb C_T).
\]
The diagram from constants to holomorphic functions induces the
specific sheaf map
\[
 \mathcal V\longrightarrow Q,
 \quad[(a_1,\ldots,a_r)]\longmapsto[(a_1,\ldots,a_r)] .
\tag{A.RT.9}\label{cdp:eq:A.RT.9}
\]
At a point $s\in S$, the $r$ entries refer to the connected
normalization germs above $s$. A tuple of branchwise constants
that comes from an element of $\OO_{S,s}$ has equal values on
all these germs, because evaluation at their common image point
is one scalar. It is therefore a diagonal tuple. This proves
that \eqref{cdp:eq:A.RT.9} is injective on every stalk. No $\OO_S$-module
structure on $\mathcal V$ is asserted or needed.

A finite map has zero positive-degree direct images of the
constant sheaf, by proper base change on its finite fibres.
Since $H^1(T;\mathbb C)=0$ and the map on constants is an
isomorphism, the connecting map from
$0\to\mathbb C_S\to\eta_*\mathbb C_T\to\mathcal V\to0$
identifies $H^0(S,\mathcal V)$ with $H^1(S;\mathbb C)$.
Naturality of the connecting maps gives the commutative square
\[
\begin{array}{ccc}
 H^0(S,\mathcal V)&\xrightarrow{\ \simeq\ }&H^1(S;\mathbb C)\\
 \big\downarrow&&\big\downarrow\\
 H^0(S,Q)&\xrightarrow{\ \simeq\ }&H^1(S,\OO_S).
\end{array}
\tag{A.RT.10}\label{cdp:eq:A.RT.10}
\]
Its left arrow is injective by \eqref{cdp:eq:A.RT.9}, proving \eqref{cdp:eq:A.RT.2}.

\ref{cdp:st:A.5.3} gives the effective integral divisor $D$ and the exact
line-bundle relation
$K_Z\otimes\OO_Z(D)\simeq(\eta\pi)^*\omega_S$.
The torsion line on the right is trivial on the smooth rational
surface: its Picard group is torsion-free and has zero identity
component. Hence $\OO_Z(-D)\simeq K_Z$. The divisor sequence
$0\to K_Z\to\OO_Z\to\OO_D\to0$, rationality of $Z$ and
Serre duality prove \eqref{cdp:eq:A.RT.3}. In detail the first positive
cohomology of $K_Z$ is zero, and $H^2(K_Z)=\mathbb C$;
the positive cohomology of $\OO_Z$ is zero. Thus $D$ is
connected. Restriction $\OO_D\to\OO_{D_{\rm red}}$ is
surjective on $H^1$, because the kernel is supported in
dimension one. Consequently $p_a(D_{\rm red})\le1$.

We identify its support, rather than discarding exceptional
components. On a minimal resolution every exceptional
irreducible curve $E$ has $K_Z\cdot E\ge0$. Indeed
$E^2<0$, and adjunction gives
$K_Z\cdot E=2p_a(E)-2-E^2$; a smooth rational
$(-1)$-curve is excluded by minimality. If $E$ has coefficient
zero in $D$, then
$0=(K_Z+D)\cdot E$ forces $K_Z\cdot E=D\cdot E=0$.
Thus $E$ is disjoint from $D_{\rm red}$.
Every exceptional fibre over the normalization conductor meets
its strict transform. Connectedness of that fibre then forces
all its exceptional curves to have positive coefficient in $D$.
Conversely, a positive-coefficient exceptional curve over a point
outside the conductor would lie in a component of $D_{\rm red}$
disjoint from every strict conductor curve. The connectedness
just proved, and the nonempty strict conductor, exclude this.
This proves the stated support equality and the zero-coefficient
assertion. It also proves connectedness of $\Gamma_{\rm red}$.

For completeness the bound on its arithmetic genus follows
directly from \eqref{cdp:eq:A.RT.8}. One has
$h^1(\Gamma,\omega_S|_\Gamma)=h^0(Q)=1+e$,
and restriction to $\Gamma_{\rm red}$ is surjective on $H^1$.
For $e=1$ the line is trivial, giving
$p_a(\Gamma_{\rm red})\le2$. For $e=0$ its $H^1$ has
dimension at most one, and degree-zero Riemann--Roch gives
$p_a(\Gamma_{\rm red})=1-h^0(\omega_S|_{\Gamma_{\rm red}})
+h^1(\omega_S|_{\Gamma_{\rm red}})\le2$.
Finally normalization and resolution are isomorphisms off the
conductor and the finitely many exceptional fibres.
Stratifying by these disjoint closed sets gives \eqref{cdp:eq:A.RT.4} by
additivity of the compactly supported Euler characteristic.
At a point outside the conductor, $T\to S$ is an
isomorphism and $S$ is normal. A normal complex surface
has isolated singularities. The reduced local equation
therefore has an isolated critical point whenever that
point is singular. Since the full fibre has equation
$u h^m$, the isolated clause of \ref{cdp:st:A.5.17} excludes it.
This proves the final assertion and eliminates the displayed
sum by an actual local comparison with the nearby torus.

We prove the asserted Euler bound for the conductor image,
including its possible positive-genus component. Put
$C=\Gamma_{\rm red}$. The normalization exact sequence for
the connected reduced curve $D_{\rm red}$ shows that the sum
of its normalized component genera is at most
$p_a(D_{\rm red})\le1$. Any component with genus one is
smooth, is unique, and meets the remaining rational curves
only in trees; a second attachment path joining distinct
points of that elliptic component would increase the
arithmetic genus to at least two. Each component of $C$ is
the finite image of a strict conductor component of
$D_{\rm red}$. Riemann--Hurwitz therefore shows that its
normalization is rational or elliptic, with at most one
elliptic-normalized component.

Write $n_0$ for the number of rational-normalized components
of $C$, $n_1\in\{0,1\}$ for the number of its
elliptic-normalized components, and $r_p,\delta_p$ for the
number of branches and the delta invariant at a singular
point. The exact coherent and topological normalization maps
give
\[
 \chi(\OO_C)=n_0-\sum_p\delta_p,
 \qquad e(C)=2n_0-\sum_p(r_p-1).
\]
For each reduced curve germ, independent constants on its
normalized branches give $\delta_p\ge r_p-1$.
The already proved arithmetic-genus bound and connectedness
give $\chi(\OO_C)\ge-1$, hence $e(C)\ge n_0-1$.
If $n_0\ge1$ this is nonnegative. If $n_0=0$, there is
one elliptic-normalized component and
$\sum\delta_p\le1$. A negative value could therefore only
be $-1$, from one ordinary node. Indeed its curve ring has
colength one in the direct sum of its two normalization
rings, and is contained in the subring of pairs with equal
constant value; equality of colengths makes these rings
equal. The two points over that node are distinct on the
elliptic normalization.

Here the hypersurface embedding excludes this last case.
At least one point $t\in T$ over the node receives both
local conductor branches. Otherwise partition the local
irreducible hypersurface branches according to which of
the two conductor branches they receive. Both parts are
nonempty. Two surface branches from opposite parts meet
in a curve, since their distinct equations in the regular
three-dimensional ambient local ring generate an ideal
of height two. That intersection curve belongs to the
conductor image, so is one of its two branches. Both
surface branches receive it, contradicting the partition.
For a locally irreducible hypersurface there is only one
normalization point and the assertion is immediate.

The connected resolution fibre $\pi^{-1}(t)$ lies in
$D_{\rm red}$ by the support equality proved above.
The strict conductor curves over the two branches are
elliptic, because a rational curve cannot map
nonconstantly to an elliptic curve. There is at most one
such curve on $D_{\rm red}$. If the two strict curves
were different, they would already give total genus two.
If they coincide, they meet $\pi^{-1}(t)$ at distinct
points: the map of that normal elliptic curve factors
through the normalization of its image, so one point
cannot map to both points over the node. The connected
resolution fibre then supplies an additional path joining
those points, outside the elliptic curve. Its components
are rational, since $T$ has rational singularities.
This produces a graph cycle in addition to the elliptic
genus, again forcing $p_a(D_{\rm red})\ge2$.
Both alternatives are impossible. Thus $e(C)\ge0$.
\end{proof}

We record the exact minimal-model calculation that follows from
\ref{cdp:st:A.5.16}. This is a reduction of the outstanding Euler calculation;
the configurations with negative complement Euler number are
retained in the calculation.
Choose any sequence of contractions of smooth rational
$(-1)$-curves from $Z$ to a minimal rational surface $Z_0$,
and let $D_0$ be the pushforward of $D$. At one contraction
$\sigma:Z_i\to Z_{i-1}$, let $E$ be the exceptional curve,
$a$ its coefficient in $D_i$, and $r$ the multiplicity of
$D_{i-1}$ at the centre. The actual divisor comparison is
\[
 K_{Z_i}+D_i
   =\sigma^*(K_{Z_{i-1}}+D_{i-1})+(1+a-r)E.
 \tag{A.RT.11}\label{cdp:eq:A.RT.11}
\]
Both canonical-plus-boundary line bundles are trivial, so
intersection with $E$ gives $a=r-1$. In particular the centre
lies on $D_{i-1}$, and effectivity is retained at every stage.
If $r=1$, the centre is a smooth point of a reduced boundary
component, $a=0$, and the strict boundary is isomorphic to the
old boundary. The blowup increases its complement Euler number
by one. If $r\ge2$, the new reduced boundary is exactly the
full inverse image of the old reduced boundary. Both the
surface and the boundary replace the centre by $\mathbf P^1$,
so the complement Euler number is unchanged. Thus
\[
 e(Z\setminus D_{\rm red})
 =e(Z_0\setminus (D_0)_{\rm red})+N_0,
 \qquad N_0\ge0,
 \tag{A.RT.12}\label{cdp:eq:A.RT.12}
\]
where $N_0$ is the number of multiplicity-one centres in this
specified sequence of reverse blowups. Formula \eqref{cdp:eq:A.RT.12} concerns
the resolution complement. For the actual torus smoothing,
the isolated clause of \ref{cdp:st:A.5.17} has already removed the
off-conductor exceptional correction in \eqref{cdp:eq:A.RT.4}. This local
argument is needed: on an abstract minimal $\mathbb F_2$,
its exceptional $(-2)$-section can be disjoint from an
anticanonical divisor without having arisen from one of the
$N_0$ blowups.

On $\mathbf P^2$, the divisor $D_0$ is an effective cubic.
Its possible reduced supports and their complement Euler
numbers are
\[
\begin{array}{c|c}
 \text{reduced support}&e(\mathbf P^2\setminus(D_0)_{\rm red})\\ \hline
 \text{smooth cubic, nodal cubic, cuspidal cubic}&3,\ 2,\ 1\\
 \text{line and conic, transverse or tangent}&1,\ 0\\
 \text{three lines forming a triangle, or concurrent}&0,\ -1\\
 \text{two distinct lines, or one line}&0,\ 1.
\end{array}
\tag{A.RT.13}\label{cdp:eq:A.RT.13}
\]
For example the union of three concurrent lines has Euler
number $6-3+1=4$, whereas $e(\mathbf P^2)=3$.
The last row retains respectively a double line plus a line,
and a triple line. These are all cubic possibilities: an
irreducible singular cubic has arithmetic genus one and
therefore exactly a node or a cusp; the other degrees partition
three as $2+1$ or $1+1+1$.

On $\mathbb F_n$, with $n=0$ or $n\ge2$, write $C_0^2=-n$
for a minimal section and $F$ for a ruling fibre. Then
\[
 D_0\sim 2C_0+(n+2)F.
 \tag{A.RT.14}\label{cdp:eq:A.RT.14}
\]
Its horizontal part has degree two over the ruling base.
If that part consists of two distinct sections $C_1+C_2$,
write $V$ for the vertical part, $d$ for its total degree and
$k=C_1\cdot C_2$. Adjunction on either section gives
$k+d=2$. Let $q$ be the number of distinct intersection
points of $C_1$ and $C_2$, and let $r$ count those points
whose ruling fibre occurs in $V_{\rm red}$. A vertical fibre
meets the two sections in two points unless they coincide at
that fibre. Additivity therefore gives the exact formula
\[
 e((D_0)_{\rm red})=4-q+r,
 \qquad e(\mathbb F_n\setminus(D_0)_{\rm red})=q-r\ge0.
 \tag{A.RT.15}\label{cdp:eq:A.RT.15}
\]
This includes two sections and a fibre through their common
point: that union has Euler number four and complement Euler
number zero.

If the horizontal part is an irreducible bisection $C$, the
same adjunction calculation gives $p_a(C)=1-d$, so $d=0$
or $1$. For $d=0$, its Euler number is at most two, and the
complement Euler number is at least two. For $d=1$, the
bisection is a smooth rational curve, the vertical part is one
reduced fibre, and their set-theoretic intersection has one or
two points. The complement Euler number is respectively one
or two.

The remaining horizontal part is $2C$, where $C$ is a section.
Write $C\sim C_0+kF$. The vertical divisor has total degree
$d=n+2-2k$, and, with $v$ its number of distinct fibres,
\[
 e((D_0)_{\rm red})=2+v,
 \qquad e(\mathbb F_n\setminus(D_0)_{\rm red})=2-v.
 \tag{A.RT.16}\label{cdp:eq:A.RT.16}
\]
If $C\ne C_0$, positivity of intersection with $C_0$ gives
$k\ge n$; hence $d\le2-n$ and $v\le2$.
For $n=0$ one also has $d\le2$. Consequently the only
negative values on a minimal ruled surface occur in the
explicit configurations
\[
 Z_0=\mathbb F_n,\quad n\ge2,\qquad
 D_0=2C_0+\sum_{j=1}^{v}a_jF_j,\qquad
 a_j\ge1,\quad \sum_j a_j=n+2,\quad v\ge3.
 \tag{A.RT.17}\label{cdp:eq:A.RT.17}
\]
The fibres $F_j$ here are distinct. Equations
\eqref{cdp:eq:A.RT.11}--\eqref{cdp:eq:A.RT.17}, together with \eqref{cdp:eq:A.RT.4}, isolate the
remaining contributions exactly; an arbitrary identification
of their conductor branches has not been substituted for an
actual hypersurface normalization map.

There is a further fundamental-group restriction on precisely
these negative minimal-model configurations. Their reduced
supports are unions of rational curves with simply connected
underlying spaces: three concurrent lines on $\mathbf P^2$,
or the section with its attached ruling fibres in \eqref{cdp:eq:A.RT.17}.
Every reverse blowup in \eqref{cdp:eq:A.RT.11} preserves this property.
A multiplicity-one blowup leaves the support isomorphic to
the preceding one. A positive-coefficient blowup at a smooth
point adds a rational leaf, at an ordinary double point
subdivides the incidence edge, and at the initial concurrent
point replaces that vertex by an exceptional $\mathbf P^1$
with three attached branches. Subsequent supports have only
the resulting ordinary double points. None of these explicit
operations introduces a topological loop. Consequently
$D_{\rm red}$ is simply connected throughout these cases.

The actual map $h=\eta\pi$ expresses $S$ as the topological
pushout of
\[
 D_{\rm red}\hookrightarrow Z,
 \qquad D_{\rm red}\xrightarrow{\ h|_{D_{\rm red}}\ }
                           \Gamma_{\rm red}.
 \tag{A.RT.18}\label{cdp:eq:A.RT.18}
\]
Indeed $h$ is an isomorphism off these closed sets and the
whole preimage of $\Gamma_{\rm red}$ is $D_{\rm red}$.
The continuous map from this pushout to $S$ is a bijection
from a compact space to a Hausdorff space, and hence is a
homeomorphism. A triangulation of the pair
$(Z,D_{\rm red})$ makes its inclusion a cofibration.
The double mapping-cylinder model for \eqref{cdp:eq:A.RT.18} therefore
has the same homotopy type as the pushout, and van Kampen
gives
\[
 \pi_1(S)\simeq
 \pi_1(Z)*_{\pi_1(D_{\rm red})}\pi_1(\Gamma_{\rm red})
 \simeq\pi_1(\Gamma_{\rm red}).
 \tag{A.RT.19}\label{cdp:eq:A.RT.19}
\]
All normalizations of the image components are $\mathbf P^1$,
by Riemann--Hurwitz. Gluing their finitely many points over
singularities identifies the topological curve with a graph
of simply connected spheres. Its fundamental group is a
free group, of rank $b_1(\Gamma_{\rm red})$.

That free group is a quotient of a metabelian group. To see
the exact comparison, triangulate $(Y,S)$ and retract a small
regular neighbourhood $N$ onto $S$, as in \ref{cdp:st:A.5.11}. Properness
puts a sufficiently small full fibre tube $U=f^{-1}(\Delta)$
inside $N$. Its retraction is surjective on fundamental
groups because it restricts to the identity on $S$.
General position makes
$\pi_1(U\setminus S)\to\pi_1(U)$ surjective, and the
torus bundle over $\Delta^*$ has
$\pi_1(U\setminus S)=\mathbb Z^4\rtimes\mathbb Z$.
This last group has abelian commutator subgroup, so its
quotient \eqref{cdp:eq:A.RT.19} is metabelian. A free group of rank at
least two is not metabelian: it maps onto $S_4$, whose
commutator group $A_4$ is not abelian. We obtain
\[
 \pi_1(S)\simeq\pi_1(\Gamma_{\rm red})\in\{1,\mathbb Z\},
 \qquad b_1(\Gamma_{\rm red})\le1.
 \tag{A.RT.20}\label{cdp:eq:A.RT.20}
\]
Under the full CDP Betti hypotheses, the nonzero class on
every support from \ref{cdp:st:D.1.9} excludes the trivial group in
\eqref{cdp:eq:A.RT.20}. This last use of the global transgression is
separate from the preceding $b_2(Y)=0$ calculations.

For these two boundary classes one can further exclude a
one-component image with first Betti number one. Such an image
has one two-branch singular point $q$; write $\alpha,\beta$
for its two normalization points there. Any two normalization
points of $S$ above $q$ receive a common local conductor branch,
because the corresponding two hypersurface branches meet in
a curve. Thus there cannot be both a point receiving only
$\alpha$ and a point receiving only $\beta$. After interchanging
the labels, every normalization point above $q$ receives
$\alpha$.

Form the incidence graph with one vertex for each strict
conductor component and one vertex for each connected
resolution fibre over a normalization point above $q$.
Edges record their intersection points, labelled by the
branch $\alpha$ or $\beta$ received there. Smooth resolution
fibres, consisting of one point, are included. The graph is
a forest: it is obtained from the tree of components and
intersection points of $D_{\rm red}$ by contracting disjoint
exceptional subtrees and deleting the other incidence edges.
Every strict conductor component maps finitely onto the sole
image component, so it has a preimage of $\beta$ and therefore
a $\beta$-edge. Every resolution-fibre vertex has an
$\alpha$-edge by the preceding paragraph. Starting at any
strict-component vertex, follow a $\beta$-edge and then an
$\alpha$-edge, and repeat. This walk never immediately reverses
an edge, because its consecutive labels differ. It is an
infinite nonbacktracking walk in a finite forest, which is
impossible. The contradiction proves the exclusion.

Since every image component is rational-normalized, write
$n_\Gamma$ for their number and $\beta_\Gamma\le1$ for the
first Betti number. Their exact topological normalization
formula now yields, on these negative boundary classes,
\[
 e(\Gamma_{\rm red})=n_\Gamma+1-\beta_\Gamma\ge2.
 \tag{A.RT.21}\label{cdp:eq:A.RT.21}
\]
For one component the excluded case is the only value below
two; for at least two components the inequality follows
directly from $\beta_\Gamma\le1$.


